\medskip

\textbf{Partie A}

\medskip

On considère deux cubes ABCDEFGH et BKLCFJMG positionnés comme sur la figure suivante:

\begin{center}
\psset{unit=1cm,arrowsize=2pt 3}
\begin{pspicture}(10.2,7.8)
\pspolygon(0.2,2.5)(5.8,0.3)(5.8,4.1)(0.2,6.3)
\psline(3,1.4)(3,5.2)(6.8,6.2)
\psline(5.8,0.3)(9.6,1.3)(9.6,5.1)(5.8,4.1)
\psline(6.8,6.2)(4,7.3)(0.2,6.3)
\psline(9.6,5.1)(6.8,6.2)
\psline[linestyle=dashed](4,7.3)(4,3.5)(9.6,1.3)
\pspolygon[linestyle=dashed](3,1.4)(6.8,2.4)(6.8,6.2)
\psline[linestyle=dashed](0.2,2.5)(4,3.5)

\psline[linewidth=1.25pt]{->}(0.2,2.5)(3,1.4)
\psline[linewidth=1.25pt,linestyle=dashed]{->}(0.2,2.5)(4,3.5)
\psline[linewidth=1.25pt]{->}(0.2,2.5)(0.2,6.3)
\psline(3,1.4)(1.6,5.75)(6.8,6.2)
\uput[dl](0.2,2.5){A} \uput[dl](3,1.4){B} \uput[ur](6.8,2.4){C} \uput[ul](4,3.5){D}
\uput[l](0.2,6.3){E} \uput[dl](3,5.2){F} \uput[ur](6.8,6.2){G} \uput[u](4,7.3){H}
\uput[u](5.8,4.1){J} \uput[d](5.8,0.3){K} \uput[r](9.6,1.3){L} \uput[ur](9.6,5.1){M}
\uput[u](1.6,5.75){I}
\end{pspicture}
\end{center}

Le point I est le milieu de [EF].

Dans toute la suite de l'exercice, on se place dans le repère orthonormé $\left(\text{A}~;~\vect{\text{AB}}~;~\vect{\text{AD}}~;~\vect{\text{AE}}\right)$.

Ainsi, par exemple, les points F, G et J ont pour coordonnées 

\begin{center}F(1~;~0~;~1),\qquad  G(1~;~1~;~1)\quad et\quad  J(2~;~0~;~1).\end{center}

\medskip

\begin{enumerate}
\item Montrer que le volume du tétraèdre FIGB est égal à $\dfrac{1}{12}$ d'unité de volume. 

On rappelle que le volume $V$ d'un tétraèdre est donné par la formule :

\begin{center}$V = \dfrac13 \times \text{aire d'une base} \times  \text{hauteur correspondante}$.\end{center}
\item Déterminer les coordonnées du point I.
\item Montrer que le vecteur $\vect{\text{DJ}}$ un vecteur normal au plan (BIG).
\item Montrer qu'une équation cartésienne du plan (BIG) est $2x - y + z - 2 = 0$.
\item Déterminer une représentation paramétrique de la droite $d$, orthogonale à (BIG) et passant par F.
\item
	\begin{enumerate}
		\item La droite $d$ coupe le plan (BIG) au point L.
		
Montrer que les coordonnées du point L sont $\left(\dfrac23~;~\dfrac16~;~\dfrac56\right)$.
		\item Calculer la longueur FL.
		\item Déduire des questions précédentes l'aire du triangle IGB.
	\end{enumerate}
\end{enumerate}

\medskip
